
\begingroup
\renewcommand{\arraystretch}{1.25}
\setlength{\LTleft}{0pt}
\setlength{\LTright}{0pt}

{\small
\begin{longtable}{@{}>{\raggedright\arraybackslash}p{0.185\textwidth}@{\hspace{6pt}}>{\raggedright\arraybackslash}p{0.175\textwidth}@{\hspace{6pt}}p{0.075\textwidth}@{\hspace{6pt}}>{\raggedright\arraybackslash}p{0.475\textwidth}@{}}
\toprule
\textbf{Notation} & \textbf{Name} & \textbf{Sec.} & \textbf{Meaning} \\
\midrule
\endfirsthead
\multicolumn{4}{@{}l}{\textit{Table of notations (continued)}}\\
\toprule
\textbf{Notation} & \textbf{Name} & \textbf{Sec.} & \textbf{Meaning} \\
\midrule
\endhead
\midrule
\multicolumn{4}{r@{}}{\textit{continued on next page}}\\
\endfoot
\bottomrule
\endlastfoot


$\T$ &
transversal gate &
all &
Transversal logical gate being gauged, with $\T^p = \id$. \\

$q, p$ & qudit dimensions & all & the original code has qudits with dimension $q$, and the added gauge qudits have dimension $p$. \\

$\physq$ &
matter-qudit set &
all &
The set of matter qudits of the input/original code. \\

$v\in \vertexV$ &
vertex qudits &
\ref{sec:gauging-procedure} &
Vertex qudits are ancillas that measure vertex stabilizers. They form vertices of the gauging graph $\graphG$. Additionally, a vertex $v$ indexes a set of matter qudits $\vertexset{v}$ and its vertex stabilizer $\av$\\

$\Sigma:\physq\to\vertexV$ &
vertex-set map &
\ref{sec:gauging-procedure} &
Assignment of each matter qudit to its vertex set. Vertex sets are constant size and respect the structure of the transversal gate $\T$. \\

$\vertexset{v}$ &
vertex set &
\ref{sec:gauging-procedure} &
The block of matter qudits assigned to vertex $v$. \\

$\T_v$ &
on-site symmetry &
\ref{sec:gauging-procedure} &
Action of the transversal gate $\T$ on vertex set $\vertexset{v}$. \\

$\Agrp$ &
symmetry group &
\ref{sec:space-distance} &
A finite Abelian group $\Agrp$. When $\Agrp$ is the cyclic group of order $p$, $\T^p=\id$. \\

$\widehat{\Agrp}$ &
representation group &
\ref{sec:space-distance} &
Additively written group of irreducible representations (characters) of $\Agrp$; the coefficient group of all representation-label chains. \\

$\T^{\bm a}$,\ $\T^g=\T^{g\id}$ &
symmetry operators &
\ref{sec:space-distance} &
$\T^{\bm a}=\prod_{v\in \vertexV}\T_v^{a_v}$ for a vertex labeling $\bm a\in \Agrp[\vertexV]$; the constant labeling where $(g\id)_v = g\ \forall v\in \vertexV$ gives the global symmetry $\T^g$. \\


$\bm\chi\in \widehat{\Agrp}[\vertexV], \omega^{\bm\chi}(\bm a)$  &
eigenoperator label on each vertex &
\ref{sec:space-distance} &
$\bm\chi$ is a labeling of vertices by $\widehat{\Agrp}$ coefficients.  $\omega^{\bm\chi}(\bm a) := \prod_v \omega^{\chi_v}(a_v)$. \\

$\Xp^a,\Zp^\chi$ &
gauge Pauli operators &
\ref{sec:gauging-procedure} &
Generalized Pauli (shift/clock) operators on $|\Agrp|$-dimensional gauge qudits; $\Xp^a\ket{x}=\ket{x+a}$, $\Zp^\chi\ket{x}=\omega^\chi(x)\ket{x}$, with strings $\Xp^{\bm a}$ and $\Zp^{\bm\alpha}$. \\

$\av^g$,\ $A(\bm a)$ &
vertex stabilizer, product of vertex stabilizers according to $\bm a\in \Agrp[\vertexV]$&
\ref{sec:gauging-procedure} & $\av^g=\T_v^g\otimes\Xp^{\delta_1 g\bm e_v}$; general finite-Abelian form $A(\bm a):=\prod_{v\in\vertexV}\av^{a_v}=\T^{\bm a}\otimes\Xp^{\delta_1\bm a}$.  For fixed $g$, the product over all vertices recovers the symmetry $\T^g$. \\

$\bp$,\ $\gaugedB_f(\chi)$ &
cycle stabilizer &
\ref{sec:gauging-procedure} &
Cycle stabilizer $\bp=\prod_{e\in\partial_2 f}\cZ_e^{\eta_f(e)}$ for a face $f\in\faceF$, whose boundary is a cycle: $\partial_1\partial_2 f=0$; we also write $\gaugedB_p(\chi)=\Zp^{\chi\,\partial_2 f}$. \\




$\graphG_\bullet, \graphG^\bullet$ &
Gauging graph chain and cochain complex. &
\ref{sec:space-distance} &
Chains are labeled by $\widehat{\Agrp}$ coefficients, cochains are labeled by $\Agrp$ coefficients. $\graphG^\bullet = \Agrp[\vertexV] \xrightarrow{\delta_{1}} 
    \Agrp[\edgeE] \xrightarrow{\delta_2} 
    \Agrp[\faceF]$, $\graphG_\bullet = \widehat{\Agrp}[\vertexV] \xleftarrow{\partial_{1}}
    \widehat{\Agrp}[\edgeE] \xleftarrow{\partial_2} 
    \widehat{\Agrp}[\faceF] $\\

$\mathcal H_{\mathrm m}$, $\mathcal H_{\mathrm g}$ &
matter and gauge Hilbert spaces &
\ref{sec:space-distance} &
$\mathcal H_{\mathrm m}$ is the Hilbert space of the original matter qudits, and
$\mathcal H_{\mathrm g}:=\bigotimes_{e\in\edgeE}\mathbb C[\Agrp]$
is the Hilbert space of the gauge qudits. \\

$\cC$,\ $\cC'$ &
base code or original code, gauged code. &
\ref{sec:gauging-procedure} &
The input stabilizer code $\cC$ with stabilizer group $\mathcal S=\langle S_i\rangle$, and the deformed code $\cC'=\langle\gaugedS_i,\av,\bp\rangle$ produced by gauging.  \\
$\cC_\mathrm{sym.}$ &
Symmetric codespace &
\ref{sec:space-distance} &
symmetric codespace is defined $\cC_\mathrm{sym.}:=\cC\cap \mathcal{H}_\mathrm{sym.}$.  \\

$\Gamma$ &
gauging map &
\ref{sec:space-distance} &
$\Gamma:\mathcal H_{\rm sym.}\xrightarrow{}\mathcal H_{\rm m}\otimes \mathcal{H}_{\rm g}$.\\


$\Gamma[M]$ & gauging map on matter operator $M$ & \ref{sec:space-distance} &
Sends a matter operator to its decorated counterpart. When $M$ is a symmetrized stabilizer, it maps it to its decorated version $\overline{\gaugedS}=\Gamma[\bar{S}]$. \\


$\aS{i}$, $\nS{i}$ &
Abelian and non-Abelian stabilizers &
\ref{sec:spacetime-distance} &
$\mathbb{Z}_2$ gauged code stabilizers can be partitioned into non-Abelian stabilizers $\nS{}$ that always commute, while Abelian stabilizers satisfy $\ComGrp{\aS{i}}{\aS{j}} \in \langle \nS{} \rangle$.\\



$\mathcal H_{\mathrm{sym.}}$ &
symmetric subspace &
\ref{sec:space-distance} &
$\mathcal H_{\mathrm{sym.}}
:=\{\ket{\psi}\in\mathcal H_{\mathrm m}:
\T^g\ket{\psi}=\ket{\psi}\ \forall g\in\Agrp\}$. \\

$\ComGrp{A}{B}$ &
group commutator &
\ref{sec:gauging-procedure} &
Group/multiplicative commutator $\ComGrp{A}{B}=ABA^\dagger B^\dagger$. \\


$\supp{O}$, $\operatorname{Wt}(O)$ &
support and weight &
\ref{sec:space-distance} &
$\supp{O}$ is the set of matter and gauge qudits on which $O$ acts nontrivially, and
$\operatorname{Wt}(O):=|\supp{O}|$. 
\\





$d$;\ $d_{\rm g}$ &
original and gauged code distance &
\ref{sec:spacetime-distance} &
The minimum weights of operators with nonscalar projected action on $\cC$ and $\cC'$ respectively. \\


$\Pi_{\cC}$, $\Pi_{\cC_\mathrm{sym.}}$, $\Pi_{\cC'}$ & code projectors & \ref{sec:space-distance} & Projectors onto the original, symmetric, and gauged codespace.\\

$M_{\bm\chi}$ & matter eigenoperator & \ref{sec:space-distance} & The component of a matter operator $M$ that satisfies $\T^{\bm a} M_{\bm\chi} (\T^{\bm a})^\dagger= \omega^{\bm\chi}(\bm a) M_{\bm\chi}.$\\

$\sum_{\bm\chi}M_{\bm\chi}$ & Eigenoperator decomposition & \ref{sec:space-distance} & Every matter operator $M$ can be written as $M=\sum_{\bm\chi\in \widehat{\Agrp}[\vertexV]} M_{\bm\chi}$\\


$\E{i}{j}\in E$ &
gauge-qubit edge &
\ref{sec:gauging-procedure} &
Gauge qubit on the edge adjacent to vertices $i, j$. \\



$\fa$, $\fn$ &
Abelian and non-Abelian faults &
\ref{sec:spacetime-distance} &
Elementary space faults on a single vertex set: $\fa$ commutes with all non-Abelian stabilizers \\

$\graphG=(\vertexV, \edgeE)= (\vertexV,\edgeE,\faceF)$ &
gauging graph &
\ref{sec:gauging-procedure} &
Graph whose vertices host the vertex ancillas (one per vertex set) and whose edges host the gauge qudits; required to have $O(1)$ degree and a sparse cycle basis. \\
 



$S$, $\gaugedS$ &
ungauged, gauged stabilizer &
\ref{sec:gauging-procedure} &
Stabilizer of the original code and its gauged counterpart. \\

$\sym$, $\overline{\gaugedS}$ &
ungauged, gauged symmetrized stabilizer &
\ref{sec:gauging-procedure} &
$\sym = \frac{1}{p}\sum_{a=0}^{p-1}\T^{a}S(\T^{a})^\dagger$, the normalized average over the symmetry group. $\overline{\gaugedS}$ denotes its gauged counterpart. \\



$h_{\Agrp}(\graphG)$ &
generalized Cheeger constant &
\ref{sec:space-distance} &
(Cheeger) expansion constant of the gauging graph.\\

$\nu$ & constant denoting the maximum number of qudits in any vertex set. & \ref{sec:space-distance} & $\nu :=\max_{v\in \vertexV}|\vertexset{v}|$\\


$\mu_v$ &
 &
\ref{sec:gauging-procedure} &
Integer defined by $\ComGrp{S_v}{\T_v}=i^{\mu_v}(N_S)_v$, constrained by $\sum_v\mu_v=0\pmod 4$. \\


$N_S$, $(N_S)_v$, $\mathbb{N}_S$ &
 &
\ref{sec:gauging-procedure} &
$N_S=\ComGrp{S}{\T}$, its restriction to vertex set $v$, and its gauged form $\mathbb{N}_S=\Circo_2[N_S]$. \\

$\Pi_{\rm sym.}$ &
symmetry projector &
\ref{sec:space-distance} &
$\Pi_{\rm sym.}=\frac{1}{|\Agrp|}\sum_{g\in \Agrp}\T^g$ \\


$\mathsf{S}$ &
generalized phase gate &
\ref{sec:gauging-procedure} &
$\mathsf{S}_P = \frac{1 + \prod_{e \in P} Z_e}{2} + i \frac{1 - \prod_{e \in P} Z_e}{2}$. \\


\end{longtable}
}
\endgroup
